% sec06 — The heat-kernel approach (writer: Writer_HeatKernel)
% Pure body LaTeX: one \section, no preamble, no \newcommand.
% Source: extract-06-heat-kernel.md; glyphs per NOTATION.md.

\section{The heat-kernel approach}\label{sec:heatkernel}

\subsection{Motivation: a semigroup-based second derivation}

The one-loop anomalous Ward identity of Section~\ref{sec:superspace}, with its
universal coefficient $\frac{\h g^2}{16\pi^2}$, is derived perturbatively: a
Schwinger--Dyson identity plus the cutting-failure mechanism of dimensional
reduction, graph by graph. A second, independent and non-perturbative
derivation of the same coefficient is highly desirable. The heat-kernel
strategy would regulate the one-loop fluctuation determinant of the full
$\mathcal N=4$ Hessian complex by the semigroup $e^{-sH}$ of a single elliptic
operator $H$ on the graded fluctuation bundle, and would extract the anomaly
from a short-time Seeley--DeWitt coefficient of that semigroup
\cite{Vassilevich}. Such a derivation would
(i)~cross-check the perturbative Schwinger-cut result without enumerating
graphs;
(ii)~make the finiteness and locality of the anomaly manifest from short-time
heat-kernel asymptotics; and
(iii)~connect directly to the evanescent structure of dimensional reduction
through the same master integral $\frac{1}{32\pi^2}$ \cite{SiegelDRED}.

This section reports the exact status of that programme. The scalar auxiliary
identities on which any such derivation must rest are established and
machine-verified (Subsection~\ref{subsec:hk:lemmas}). The full derivation,
however, is blocked at the regulator level: no single elliptic generator on
the full $\mathcal N=4$ fluctuation complex has been constructed, and the
anomaly coefficient used elsewhere in this paper is \emph{not} obtained from
this route. Subsection~\ref{subsec:hk:blocker} states the obstruction
precisely, and Subsection~\ref{subsec:hk:closure} lists what would close the
route.

\subsection{Established scalar and semigroup lemmas}\label{subsec:hk:lemmas}

Every identity in this subsection is exact; each is proved by direct scalar
computation and is machine-verified in the source repository. We write
$H(g)=H_0+gH_1+g^2H_2$ for the one-loop fluctuation operator expanded in the
coupling, and $s>0$ for the proper-time parameter.

\begin{lemma}[Second-order Duhamel coefficient]\label{lem:hk:duhamel}
Let $E\to M$ be a fixed graded bundle, let $H_0:E\to E$ be an even operator,
and let $H_1,H_2$ be even endomorphism-valued differential operators. Set
\begin{equation}\label{eq:hk:Hg}
H(g)=H_0+gH_1+g^2H_2 .
\end{equation}
Then the coefficient of $g^2$ in the semigroup $e^{-sH(g)}$ is the two-term
ordered integral
\begin{equation}\label{eq:hk:duhamel}
\begin{aligned}
[g^2]\,e^{-sH(g)}
={}&-\int_0^s dt\; e^{-(s-t)H_0}\,H_2\,e^{-tH_0} \\
&+\int_{0<t_1<t_2<s} dt_1\,dt_2\;
e^{-(s-t_2)H_0}\,H_1\,e^{-(t_2-t_1)H_0}\,H_1\,e^{-t_1H_0} .
\end{aligned}
\end{equation}
\end{lemma}

\begin{remark}[The quadratic background block cannot be omitted]
The first line of \eqref{eq:hk:duhamel} is at the same background order as the
second line. Hence a heat-kernel derivation of the physical coefficient must
derive the full $H_2$ seagull, connection-square, gauge and ghost blocks; the
quadratic background block cannot be omitted in favor of two linear $H_1$
insertions.
\end{remark}

\begin{lemma}[Ordered-simplex moment]\label{lem:hk:simplex}
For integers $p,q\ge0$,
\begin{equation}\label{eq:hk:simplex}
M_{p,q}:=\int_{0<a<b<1} a^p b^q \,da\,db
=\frac{1}{(p+1)(p+q+2)} .
\end{equation}
Consequently, for $0\le k\le m$ and $0\le\ell\le n$, the tower coefficient is
\begin{equation}\label{eq:hk:tower}
T_{m,n;k,\ell}
:=\binom{m}{k}\binom{n}{\ell}\,M_{k+\ell,\,m+n-k-\ell}
=\frac{\binom{m}{k}\binom{n}{\ell}}{(m+n+2)(k+\ell+1)} .
\end{equation}
\end{lemma}

\begin{proof}
Integrating $a$ at fixed $b$ gives
$M_{p,q}=\int_0^1 db\,b^q\,\frac{b^{p+1}}{p+1}=\frac{1}{(p+1)(p+q+2)}$,
and \eqref{eq:hk:tower} follows by substitution of \eqref{eq:hk:simplex}.
\end{proof}

\begin{remark}[A second ordering supplies no automatic factor two]
\label{rem:hk:factortwo}
The moment \eqref{eq:hk:simplex} sums a single ordering of the simplex; no
factor from the reversed ordering is included. A factor of two may be added
only after an independent ordered-vertex calculation proves equality of the
two orderings, including color order and Koszul sign.
\end{remark}

\begin{lemma}[Gaussian convolution and bridge means]\label{lem:hk:gaussian}
For $s>0$ and $x\in\mathbb R^4$ let
\begin{equation}\label{eq:hk:kernel}
K_s(x)=\frac{1}{(4\pi s)^2}\,e^{-x^2/4s} .
\end{equation}
Then $K_s$ is normalized, solves the heat equation, and has diagonal value
\begin{equation}\label{eq:hk:heateq}
\int_{\mathbb R^4} d^4x\,K_s(x)=1,\qquad
\partial_s K_s=\Delta_{\mathbb R^4}K_s,\qquad
K_s(0)=\frac{1}{16\pi^2 s^2} ,
\end{equation}
and it obeys the convolution (Chapman--Kolmogorov) identity
\begin{equation}\label{eq:hk:conv}
\int_{\mathbb R^4} d^4y\,K_s(x-y)K_t(y)=K_{s+t}(x) .
\end{equation}
Moreover, for the three-segment bridge with $t_1,t_2,t_3>0$,
$T=t_1+t_2+t_3$, and Gaussian weight
$Q(x,y)=\frac{(x-w)^2}{4t_1}+\frac{(y-x)^2}{4t_2}+\frac{y^2}{4t_3}$ in one
Cartesian component, the stationary means are
\begin{equation}\label{eq:hk:bridge}
\langle x\rangle=\frac{t_2+t_3}{T}\,w,\qquad
\langle y\rangle=\frac{t_3}{T}\,w ,
\end{equation}
and the inverse Hessian of $Q$ is homogeneous of degree one under
$(t_1,t_2,t_3)\mapsto(\rho t_1,\rho t_2,\rho t_3)$, $\rho>0$.
\end{lemma}

\begin{remark}
Equation \eqref{eq:hk:bridge} proves the exact bridge mean and covariance
scaling; it does not remove the endpoint factor $K_T(w)$.
\end{remark}

\begin{lemma}[Master-integral consistency with dimensional reduction]
\label{lem:hk:dred}
Let $d=4-2\varepsilon$, and let $\mul^2$ be the evanescent component of the
loop momentum squared (Section~\ref{sec:gaugebv}). Then
\begin{equation}\label{eq:hk:dred}
\lim_{\varepsilon\to0}\mu^{2\varepsilon}
\int\frac{d^d\ell}{(2\pi)^d}\,
\frac{\mul^2}{(\ell^2+\Delta)^3}
=\frac{1}{32\pi^2} .
\end{equation}
\end{lemma}

\begin{remark}
The value \eqref{eq:hk:dred} agrees with the master integral of the accepted
derivation (Sections~\ref{sec:superspace} and \ref{sec:component}), so the
scalar heat-kernel layer is consistent with the evanescent mechanism of
dimensional reduction \cite{SiegelDRED}. Equality of this scalar integral does
not, however, determine the $\mathcal N=4$ operator pairing, the color tensor,
the orientation sign, or the anomaly multiplicity.
\end{remark}

\begin{lemma}[Vanishing short-time limit at fixed nonzero shift]
\label{lem:hk:shorttime}
For fixed $w\neq0$,
\begin{equation}\label{eq:hk:shorttime}
\lim_{s\downarrow0}\frac{1}{(4\pi s)^2}\,e^{-w^2/4s}=0 .
\end{equation}
Consequently a local derivative tower requires distributional coefficient
extraction at $w=0$; it cannot be obtained by deleting $K_s(w)$ before taking
the limit.
\end{lemma}

\subsection{The obstruction: the regulator proposal fails}
\label{subsec:hk:blocker}

The proposed continuation defines its would-be generator blockwise from a
Hessian two-form $\delta^2S$ and a field-space identification $\mathcal G$.
Its own displayed free matter blocks give
\begin{equation}\label{eq:hk:matter}
\big(\mathcal G^{-1}\delta^2S\big)^2\Big|_{E_+}
=\bar\nabla^2\nabla^2
=16\,\BoxE\,\mathcal P_+ ,
\end{equation}
whereas the declared generator is $-\BoxE\,\mathcal P_+$; the accompanying
verifier inserts a factor $-1/16$ that appears nowhere in the proposal. Both
the normalization ($16$) and the sign are wrong relative to the declared
generator. Likewise, the proposal's own vector-block weights give
\begin{equation}\label{eq:hk:vector}
\mathcal G_V^{-1}\,\delta^2S_{VV}
=\big(2g^2\big)\frac{1}{2g^2}\,\BoxE
=\BoxE ,
\end{equation}
whereas the declared generator is $-\BoxE$; the verifier adds an undeclared
Euclidean sign. Finally, the prescription applies the squared operator on the
matter and ghost rows but the unsquared operator on the vector row. This is
not a single operator functional calculus when nonzero vector--matter mixing
blocks are retained: different powers on different blocks define different
operators, so the off-diagonal mixing blocks have no defined heat evolution
under this prescription.

\begin{remark}[Status of the route: blocked]\label{rem:hk:blocked}
For the three independent reasons above, the full heat-kernel derivation of
the anomaly coefficient is currently blocked; the source repository records
the status
\texttt{BLOCKED\_HEAT\_KERNEL\_TYPED\_REGULATOR\_AND\_COEFFICIENT\_DERIVATION}.
The identities of Subsection~\ref{subsec:hk:lemmas} are admitted as verified
auxiliary lemmas only: they do not define a regulator on the full
$\mathcal N=4$ fluctuation complex, and they do not independently derive the
coefficient $\frac{\h g^2}{16\pi^2}$, which remains the result of the
Schwinger-cut derivation of Section~\ref{sec:superspace}.
\end{remark}

\subsection{What would close the route}\label{subsec:hk:closure}

The route would produce a regulator only when the following are derived
together, not declared:
\begin{itemize}
\item a \emph{single elliptic generator} on the full graded $\mathcal N=4$
fluctuation bundle---one operator, one functional calculus, with the correct
sign and normalization on every block (matter, vector, ghost), covering also
the off-diagonal vector--matter mixing blocks;
\item its \emph{dual Euler action} on the letter bilinears, consistent with
the Schwinger--Dyson/Euler reduction of Section~\ref{sec:superspace};
\item the \emph{full quadratic coefficient} $H_2$---the seagull,
connection-square, gauge and ghost blocks at order $g^2$---which by
Lemma~\ref{lem:hk:duhamel} sits at the same background order as two linear
insertions and cannot be dropped;
\item the \emph{distributional short-time limit}: coefficient extraction at
$w=0$ in the distributional sense, as required by Lemma~\ref{lem:hk:shorttime}.
\end{itemize}
Reproducing the accepted coefficient $\frac{\h g^2}{16\pi^2}$ with the correct
color tensor and the correct census of nonzero channels would in addition
require the $\mathcal N=4$ operator pairing, the orientation sign and the
anomaly multiplicity to be derived from that generator, since none of them
follow from the verified scalar identities alone. Until then, the heat-kernel
approach stands as an auxiliary consistency layer, not as a second derivation.
